%\usepackage [ISO-8859-1]{inputenc}
\documentclass[a4paper,12pt,oneside]{amsbook}
\usepackage[T1]{fontenc}
\usepackage{listings}
\usepackage{textcomp}
\usepackage{multicol,calc}
%\usepackage{euler}
%\usepackage{garamond}
\usepackage{hyperref}

\usepackage{newcent}
%\fontfamily{fxm}\selectfont
\usepackage{fancyhdr}
\pagestyle{fancy}
% Detta paket skoter grafiken
%\usepackage[dvips]{graphicx}
\usepackage[pdftex]{graphicx}

\usepackage{amsthm}
\usepackage{amsmath}
\usepackage[swedish]{babel}
% For bortkommentering
\usepackage{verbatim}
% For att slippa indentering av Sections och SubSections
\usepackage{indentfirst}

% Har stalls paragrafindentering och avstand mellan stycken och rader mm
\setlength{\parindent}{0pt}
\setlength{\parskip}{6pt}
\setlength{\baselineskip}{12pt}
% Har definieras sidan
\setlength{\voffset}{-15mm}
\setlength{\hoffset}{-5.5mm}

\setlength{\topmargin}{0mm}
\setlength{\headheight}{10mm}
\setlength{\headsep}{8mm}
\setlength{\footskip}{16mm}

\setlength{\evensidemargin}{0mm}
\setlength{\oddsidemargin}{0mm}

\setlength{\marginparsep}{0mm}
\setlength{\marginparwidth}{0mm}

\setlength{\textwidth}{170mm}
\setlength{\textheight}{230mm}
\setlength{\paperwidth}{210mm}
\setlength{\paperheight}{297mm}
\setlength{\headwidth}{170mm}
% Huvud och Fot
%\fancypagestyle{}{
\fancyhf{}\fancyhead[c]{\small{\textit{Programmeringsolympiaden Kvalificering 2025}}}
\fancyfoot[c]{\thepage}
\renewcommand{\headrulewidth}{0.4pt}
\renewcommand{\footrulewidth}{0.4pt}
\newenvironment{lista}
{\begin{itemize}
\setlength{\parindent}{0pt}
\setlength{\itemsep}{6pt}}
{\end{itemize}}
% Uppgifter
\newcounter{probnr}
\newenvironment{tal}{%
\begin{list}
%{\textbf{\arabic{section}.\arabic{probnr}}} {\usecounter{probnr}
{\textbf{\arabic{probnr}}} {\usecounter{probnr}
\setlength{\leftmargin}{0mm}
\setlength{\rightmargin}{0mm}
\setlength{\labelwidth}{-1mm}
\setlength{\labelsep}{1mm}}
\setlength{\itemsep}{6pt}
}{\end{list}}
% Definition av avsnitt: FAKTA, EXEMPEL, PASTAENDE
\newtheorem{fakta}{Fakta}
\newtheorem{exempel}{Exempel}
\newtheorem{problem}{Problem}
%
\newtheoremstyle{test}% NAME
{20pt}      % ABOVESPACE
{10pt}      % BELOWSPACE
{\sffamily} % BODYFONT
{0pt}       % INDENT
{\scshape}  % HEADFONT
{}          % HEADPUNCT
{\newline}  % HEADSPACE
{}          % CUSTOM-HEAD-SPEC

\theoremstyle{test}
\newtheorem{program}{Program}
\newcommand{\sv}[1]{\textsc{#1}}            % Sma versaler
\newcommand{\fe}[1]{\textbf{#1}}            % FET
\newcommand{\ku}[1]{\textit{#1}}            % KURSIV
\newcommand{\cu}[1]{\texttt{#1}}            % Courier
\newcommand{\sk}[1]{\texttt{#1}}            % Courier
\newcommand{\rubrik}[1]{\begin{center}\sf\huge{#1}\normalsize\rm\end{center}}
\begin{document}
%\DeclareGraphicsExtensions{.jpg,.pdf,.mps,.png,.eps}



\thispagestyle{fancy}
\lstset{basicstyle=\ttfamily,
  breaklines=true}
% ----------------------------------------------------

\begin{center}
\Huge{Programmeringsolympiaden 2025}
\end{center}
\vspace{-1.2cm} 
\specialsection*{Tävlingsregler för skolkvalet}
\begin{lista}
\item Tävlingen äger rum på av skolan bestämt datum under {\bf fyra
    timmar. Ingen förlängning ges för lunch eller raster.} Eleven ska i förväg komma överens med läraren om att använda egen dator eller en som skolan tillhandahåller. I vilket fall som helst måste eleven befinna sig i avtalad lokal på skolan. 
    %{\bf På grund av pandemin kan i undantagsfall tävling ske från annan plats godkänd av läraren, givet tillräcklig övervakning.}
\item Tävlingen består av sju uppgifter som vardera ska lösas genom ett
  datorprogram i valfritt programmeringsspråk.
\item {\bf Indata kan läsas in i programmet på valfritt sätt}, t.ex. genom att programmet för en dialog med användaren (som i körningsexemplena i uppgifterna), att de skrivs in i ett grafiskt gränssnitt eller att datafiler skickas till {\em standard input}. Kom bara överens med din lärare om hur programmet ska testas.
\item Dina lösningar kommer att testköras med förpreparerade
  indata. Varje uppgift testas normalt med 5 testfall, som vardera ger
  1 poäng om ditt program skriver ut korrekt svar inom en exekveringstidsgräns av
  {\bf 3 sekunder}. Ingen test av indata behöver göras, den följer specifikationerna 
i uppgiften.
\item Det är ofta olika begränsningar på de olika testfallen, t.ex.\ storleken på
indata eller andra inskränkningar. Detta anges i uppgiften. {\bf Observera att det kan vara
helt olika svårighetsgrad på en uppgift beroende på dessa skillnader. Det kan därför vara
lättare att få delpoäng på en uppgift som verkar svår än att få full
poäng på en uppgift som verkar lättare.} Informationen om delpoäng är
därför extremt viktig för att planera sin
tävling.
\item Rättningen utförs på samma eller likvärdig dator. Ändringar
  i källkoden tillåts ej efter tävlingen. Om programmet
  inte kan kompileras ges 0 p.\ på uppgiften. 
\item Om något av följande inträffar ger
  det {\em testfallet} 0 poäng, men programmet fortsätter testas med övriga testfall.
\begin{lista}
\item Exekveringstiden överstiger 3 sekunder \vspace{-0.2cm}
\item Exekveringsfel (run time error) \vspace{-0.2cm}
\item Fel svar \vspace{-0.2cm}
\end{lista}

\item Deltagandet är individuellt vilket bland annat innebär att inget utbyte av idéer eller 
filer får ske under tävlingen.  
\item Hjälpmedel: Valfritt skriftligt material, material som finns installerat på datorn samt material som finns tillgängligt på internet.
Det är \ku{inte} tillåtet att aktivt kommunicera på internet (t.ex. chatta eller ställa frågor till ett forum) utan endast att söka efter information.
{\bf Det är heller inte tillåtet att använda generativ AI såsom ChatGPT eller Github Copilot.}
Räknedosa är tillåten.
\item Tävlingsbidraget ska lämnas in i form av källkodsfiler (uppg1...uppg7 med passande filtillägg) som läggs på överenskommen plats. Ingen hänsyn tas till andra filer. Var noga 
med att lämna in den korrekta versionen av ditt program.  
\end{lista}

%De högst placerade i kvalet går vidare till finalen i Göteborg där landslagsplatser till BOI i Polen och IOI i Bolivia står på spel.
\begin{center}
\Large Lycka till!
\end{center}
%\end{document}

\newpage
\specialsection*{Uppgift 1 -- Triangelfabriken}

Tristian jobbar på en triangelfabrik.
Hans arbete är att klassificera olika sorters trianglar som produceras i fabriken.
Tristian ber nu dig att skriva ett program som kan ersätta honom.

Tristian kommer att ge dig 3 positiva heltal, $a$, $b$ och $c$, som är de 3 vinklarna i en triangel.

Ditt program ska göra följande:
\begin{itemize}
  \item Om triangeln är trubbig, så ska programmet skriva ut ``Trubbig Triangel''.
  \item Om triangeln är spetsig, så ska programmet skriva ut ``Spetsig Triangel''.
  \item Om triangeln är rätvinklig, så ska programmet skriva ut ``Ratvinklig Triangel''.
  (Tristian har inte lärt sig alla bokstäver i det svenska alfabetet än, så han förstår inte om du skriver ut ``Rätvinklig'' med ett ``ä''. Han förstår bara ``Ratvinklig Triangel'').
\end{itemize}

För att påminna om vad trubbiga, spetsiga och rätvinkliga trianglar är:
\begin{itemize}
  \item En triangel är trubbig om någon av dess vinklar är större än $90$ grader.
  \item En triangel är spetsig om alla dess vinklar är mindre än $90$ grader.
  \item En triangel är rätvinklig om en av dess vinklar är exakt $90$ grader.
\end{itemize}

Ditt program ska först läsa in de tre vinklarna $a$, $b$, $c$.

$a$ kommer stå på första raden.

$b$ kommer stå på andra raden.

$c$ kommer stå på tredje raden.

Det är garanterat att $(1 \leq a, b, c < 180)$ håller. Det är också garanterat att indatan är en triangel, alltså att $a + b + c = 180$.

Sen ska programmet skriva ut om triangeln som utgörs av vinklarna $a$, $b$, $c$ är spetsig, trubbig eller rätvinklig. Var noga
med att det du skriver ut är samma som det där uppe! Det vill säga, ``Trubbig Triangel'' om den är trubbig, etc.

\fe{Poängsättning:} Programmet kommer att köras på 5 testfall, vardera värt $1$ poäng.

\vspace{1cm}
\fe{Körningsexempel 1}
\begin{verbatim}
Vinkel 1 ? 100
Vinkel 2 ? 40
Vinkel 3 ? 40

Svar: Trubbig Triangel
\end{verbatim}
Se nästa sida för fler körningsexempel.

\newpage

\fe{Körningsexempel 2}
\begin{verbatim}
Vinkel 1 ? 60
Vinkel 2 ? 40
Vinkel 3 ? 80

Svar: Spetsig Triangel
\end{verbatim}

\fe{Körningsexempel 3}
\begin{verbatim}
Vinkel 1 ? 30
Vinkel 2 ? 60
Vinkel 3 ? 90

Svar: Ratvinklig Triangel
\end{verbatim}


\newpage
\specialsection*{Uppgift 2 -- Tre i rad}


%\begin{center}
%  \includegraphics[width=12cm]{sifferkrypto.pdf} \\ ~\\
%  \ku{Exempel på missuppfattning som kan ske med Alice nya chiffer.}
%\end{center}

\vspace{1cm}

Matematikern Lenore Oljer har gjort så mycket matte att hon har tröttnat på primtal. 
Hon har nu börjat studera så kallade trevliga tripletter. 
Det är väldigt lätt att skapa en trevlig triplett. 
Börja med ett heltal $a$ som är
större än noll och skapa sedan tripletten $(a,a+1,a+2)$. 
Alla tripletter som skapats på detta sätt är trevliga. 
Några exempel på trevliga tripletter är $(4,5,6)$ och $(15,16,17)$.

För att lära sig mer om trevliga tripletter vill Lenore hitta hur 
många tal som kan skrivas som produkten av
alla talen i en trevlig triplett. Hon kallar dessa tal ``treiga tal''. 
Några exempel på treiga tal är $24=2 \cdot 3 \cdot 4$
och $336 = 6 \cdot 7 \cdot 8$. Lenore vill nu veta hur många treiga tal det 
finns som är mindre än talet $N$.

Skriv ett program som läser in talet $N$ och skriver ut hur många tal som är mindre än
$N$ och är treiga.

\fe{Poängsättning:}\\
De fem testfallen som vardera ger $1$ poäng har följande begränsningar:
\begin{enumerate}
\item  $\;\;\;\;N \le 25 $    
\item  $\;\;\;\;N \le 80 $ 
\item  $\;\;\;\;N \le 1000 $   
\item  $\;\;\;\;N \le 10^5 $   
\item  $\;\;\;\;N \le 10^9 $   
\end{enumerate}


\vspace{1cm}

\fe{Körningsexempel 1}
\begin{verbatim}
N ? 75

Svar: 3
\end{verbatim}

\fe{Förklaring:} De treiga talen mindre än 75 är 6, 24 och 60.

\vspace{0.5cm}

\fe{Körningsexempel 2}
\begin{verbatim}
N ? 24

Svar: 1
\end{verbatim}

\fe{Förklaring:} Eftersom det treiga talet $24$ inte är mindre än $N$ räknar vi bara 6.

\vspace{0.5cm}

\fe{Körningsexempel 3}
\begin{verbatim}
N ? 1234

Svar: 9
\end{verbatim}



\newpage
\specialsection*{Uppgift 3 -- Meddelande}

I staden som Paula bor i finns ett stort torg. Torget kan representeras av ett 
rutnät med $n$ rader och $m$ kolumner. I vissa av rutorna finns en bokstav 
inristad. 

Enligt legenden så bildar dessa bokstäver ett meddelande om man 
går på rutorna i en viss ordning. För att det ska bli rätt måste man börja 
i det övre vänstra hörnet och sluta i det nedre högra hörnet, och man får bara gå steg åt höger 
och nedåt. Om man lyckas besöka alla bokstavsrutorna på det här viset så bildas 
det hemliga meddelandet.

Paula har beslutat sig för att en gång för alla knäcka gåtan, och hitta meddelandet.

Skriv ett program som först läser in antalet rader ($n$) och kolumner ($m$) som rutnätet har. Läs sedan
in $n$ rader, som innehåller alla rutnätets rader. För alla testfall håller det att $1 \leq n,m \leq 6$.

Programmet ska sedan skriva ut det hemliga meddelandet som står på torget.

\vspace{0.5cm}
\fe{Poängsättning:}\\
De fem testfallen som vardera ger $1$ poäng har följande begränsningar utöver de som redan nämnts:
\begin{itemize}
    \item[(1)]   $\;\;\;\;$ $n=1$    
    \item[(2-3)] $\;\;\;\;$ Antalet bokstäver är $n+m-1$ 
    \item[(4-5)] $\;\;\;\;$ Inga särskilda begränsningar.
\end{itemize}

\vspace{1cm}
\fe{Körningsexempel 1}\\
\begin{verbatim}
n ? 5
m ? 6
Rad 1 ? pa....
Rad 2 ? ......
Rad 3 ? .u.l..
Rad 4 ? .....a
Rad 5 ? ......

Svar: paula
\end{verbatim}

\begin{center}
  %\includegraphics[width=6cm]{meddelande.png} \\ ~\\
  \ku{Illustration av rutnätet i körningsexempel 1.}
\end{center}


\fe{Körningsexempel 2}\\
\begin{verbatim}
n ? 3
m ? 3
Rad 1 ? sn.
Rad 2 ? .a.
Rad 3 ? .ke

Svar: snake
\end{verbatim}


\newpage
\specialsection*{Uppgift 4 -- Klippa protein}

Din kompis Björn har nyligen påbörjat ett stort kemiexperiment, med målet att skapa ett hemligt protein. 
Alla protein beskrivs med strängar av bokstäver från \texttt{a}-\texttt{z}, exempelvis \texttt{bcaa} och \texttt{hejsan}.
Han behöver skapa ett flertal så kallade \textit{rena} proteiner. Dessa är proteiner som består av kopior av en och samma bokstav.
För att göra detta har han mikroskopiska saxar som kallas CRISPR. Med hjälp av dessa kan han ta bort upp till $k$ bokstäver
som ligger bredvid varandra. Med andra ord kan han välja ett intervall med upp till $k$ bokstäver att klippa bort.
Om ett protein delas i två delas när det klipps kommer dessa delar sedan att sättas ihop där de klipptes av. 

Eftersom han är väldigt lat vill han ha din hjälp att skriva ett program som beräknar det minsta antalet gånger han behöver klippa ett visst protein för
att det ska bli rent. Notera att proteinet \textbf{inte} får bli en tom sträng, eftersom det isåfall inte längre kan användas för kemi.

Programmet ska först läsa in Björns protein, som består av som mest $32$ bokstäver. Sedan ska det läsa in talet $k$ ($k \ge 1$).

Programmet ska sedan skriva ut det minsta antalet klippningar som krävs för att göra proteinet rent.

\vspace{0.5cm}
\fe{Poängsättning:}\\
De fem testfallen som vardera ger $1$ poäng har följande begränsningar utöver de som redan nämnts:
\begin{itemize}
    \item[(1)] $\;\;\;\;$ Strängen består endast av bokstäverna \texttt{a} och \texttt{b}.    
    \item[(2)] $\;\;\;\;k=1$ 
    \item[(3-5)] $\;\;\;\; $Inga särskilda begränsningar.
\end{itemize}

\vspace{1cm}
\fe{Körningsexempel 1}
\begin{verbatim}
Björns protein ? exempelfall
k ? 4

Svar: 3
\end{verbatim}

\fe{Förklaring:}
Ett möjligt sätt är först klippa bort bokstäverna \texttt{pe}. Vi får då strängen \texttt{exemlfall}. \\
Sedan kan vi klippa bort \texttt{fa} och får \texttt{exemlll}. \\
Till sist klippa vi bort \texttt{exem} och får \texttt{lll}. 

\vspace{3cm}

\fe{Körningsexempel 2}
\begin{verbatim}
Björns protein ? aabbabbba
k ? 3

Svar: 2
\end{verbatim}


\begin{center}
  %\includegraphics[width=12cm]{klippaprotein.png} \\ ~\\
  \ku{Illustration av hur Björn kan klippa proteinet optimalt i körningsexempel 2.}
\end{center}

\newpage
\specialsection*{Uppgift 5 -- Miniräknaren}

Det går rykten om att skolans miniräknare kan trimmas med en hemlig kod.
Du har hört att om du kan få miniräknaren att visa den hemliga kålnami-koden på displayen och sedan
trycka på \texttt{=}-tecknet, så kommer miniräknaren att låsa upp flera
stycken hemliga matematikfunktioner.

Din kompis har lyckats lista ut vad kålnami-koden är, men han berättar för dig
att du inte får gå tillväga hursomhelst för att skriva in koden i miniräknaren.
Från början visar miniräknaren $0$. Sedan får du utföra följande operationer:
\begin{itemize}
  \item Multiplicera talet som visas med talet $m$.
  \item Addera ett valfritt tal $k$ till talet som visas, där $k$ uppfyller $0 \leq k \leq m-1$.
\end{itemize}

Som tur är har din kompis även lyckats lista ut vad talet $m$ är. 
Denna komplicerade process kommer endast att fungera om du gör så få operationer
som möjligt. Därför vill du nu ta reda på hur många operationer som krävs,
givet kålnami-koden och talet $m$.

Programmet ska först läsa in heltalet $n$ ($1 \le n \le 10^9$), innehållet av kålnami-koden.
På nästa rad ska det läsa in siffran $m$ ($2 \leq m \leq 9$), som beskrivits ovan.

Programmet ska sedan skriva ut ett heltal: minsta antalet operationer som krävs för att få miniräknaren att visa
kålnami-koden $n$.

\fe{Poängsättning:}\\
De fem testfallen som vardera ger $1$ poäng har följande begränsningar utöver de som redan nämnts:
\begin{itemize}
    \item[(1)] $\;\;\;\;m=3, n \leq 10$ 
    \item[(2)] $\;\;\;\;m=2$ 
    \item[(3)] $\;\;\;\;n \leq 10^5$
    \item[(4-5)] $\;\;\;\; $Inga särskilda begränsningar.
\end{itemize}


\fe{Körningsexempel 1}
\begin{verbatim}
n ? 4
m ? 2

Svar: 3
\end{verbatim}

\fe{Förklaring:}
Ett sätt att skriva in koden med minsta antalet operationer är att göra som följande:
\begin{enumerate}
  \item $+1$
  \item $\times 2$
  \item $\times 2$
\end{enumerate}


\fe{Körningsexempel 2}
\begin{verbatim}
n ? 33
m ? 3

Svar: 8
\end{verbatim}
\fe{Förklaring:}
En optimal lösning är att utföra följande operationer:
\begin{enumerate}
  \item $+1$
  \item $\times 3$
  \item $\times 3$
  \item $+2$
  \item $\times 3$
\end{enumerate}


\newpage
\specialsection*{Uppgift 6 -- Trollkarlen Theodor}

En mindre armé med $N$ monster befinner sig precis utanför Tästerås! Nu är det upp
till Trollkarlen Theodor att besegra den och rädda staden. Varje monster har en viss
mängd liv. För att skada monster kan
Theodor avfyra magiska explosioner. Varje gång han avfyrar en explosion siktar han på ett
särskilt monster. Sen händer följande:
\begin{itemize}
  \item Alla monster förlorar $A$ liv av explosionen.
  \item Monstret som Theodor siktade på förlorar $S$ liv
  till, eftersom det befinner sig i mitten av explosionen. Alltså förlorar monstret han siktar
  på $S+A$ liv sammanlagt av just den explosionen.
\end{itemize}
Han kan fritt välja vilket monster han ska sikta på varje gång han avfyrar en explosion.
Eftersom det är jobbigt att avfyra explosioner vill han veta minsta antalet
explosioner han behöver skjuta iväg för att döda alla monster om han siktar på monster optimalt.

Skriv ett program som läser in heltalen $N$ ($1 \leq N \leq 10$), $S$ ($1 \leq S \leq 10^9$) och $A$ ($0 \leq A \leq 10^9$), 
vars betydelse beskrivs ovan. På nästa rad ska programmet läsa in $N$ heltal $h_1, h_2, \ldots, h_N$ ($1 \leq h_i \leq 10^9$), som beskriver antalet liv som vardera monster har.

Programmet ska sedan skriva ut minsta antalet explosioner som Theodor behöver avfyra för att döda
alla monster.

\fe{Poängsättning:}\\
De fem testfallen som vardera ger $1$ poäng har följande begränsningar utöver de som redan nämnts:
\begin{itemize}
    \item[(1)] $\;\;\;\; A=0$
    \item[(2-3)] $\;\;\;\; h_i \leq 100$
    \item[(4-5)] $\;\;\;\; $Inga särskilda begränsningar.
\end{itemize}


\fe{Körningsexempel 1}
\begin{verbatim}
N, S, A ? 3 2 1
Monstrenas liv ? 7 2 3

Svar: 3
\end{verbatim}
Ett optimalt sätt att avfyra explosioner är följande:
\begin{enumerate}
  \item Avfyra en explosion på monstret med 7 liv. Monstrena har nu \verb+4,1,2+ liv kvar.
  \item Avfyra en explosion på monstret med 4 liv kvar. Monstret med liv 1 dör. De resternade monstrena har \verb+1,1+ liv kvar. 
  \item Avfyra en explosion på vilket som helst av monstrena. Nu har alla monster dött och Theodor är klar.
\end{enumerate}


\fe{Körningsexempel 2}
\begin{verbatim}
N, S, A ? 4 3 2
Monstrenas liv ? 7 8 9 10

Svar: 4  
\end{verbatim}


\specialsection*{Uppgift 7 -- Theodor, Trollkarlen}


Detta problem bygger på problemet Trollkarlen Theodor.

Ondskefulla Olle, som är ytterst ansvarig för all ondska, har tröttnat på att Theodor alltid förstör allting.
Nu har Olle kommit på en ny ondskefull plan, och för att sätta den i verket måste han uppehålla Theodor.
Till sitt förfogande har han $N$ monster, och eftersom han är så enormt mäktig kan han ge dessa monster olika antal liv.
Olle vill nu räkna antalet sätt han kan ge olika mängder liv till sina olika monster, 
så att Theodor behöver avfyra exakt $K$ magiska explosioner för att döda alla monster. Enda kravet på antalet liv
för ett monster är att det ska vara större än $0$.

Precis som i problemet Trollkarlen Theodor gör Theodor först $S$ skada på monstret han siktar på, och sedan $A$ skada på alla monster (även det han siktade på).
Theodor väljer även hur han siktar på monster så att han minimerar antalet explosioner han behöver avfyra.


Skriv ett program som läser in heltalen, $N, K, S, A$ ($1 \leq N,K,S \leq 2\cdot10^5$, $0 \leq A \leq 2\cdot10^5$), vilket är de fyra talen som beskrivits ovan.

Programmet ska sedan skriva ut minsta antalet explosioner som Theodor behöver avfyra för att döda
alla monster. Eftersom detta talet kan bli väldigt stort ska det skrivas ut modulo $10^9+7$. Om du bara vill lösa
första testfallet behöver du inte tänka på modulo eftersom svaret är litet.


\fe{Poängsättning:}\\
De fem testfallen som vardera ger $1$ poäng har följande begränsningar utöver de som redan nämnts:
\begin{itemize}
    \item[(1)] $\;\;\;\; A = 0$, $N, K, S \leq 5$
    \item[(2)] $\;\;\;\; A = 0$, $S \leq 10$, $N, K \leq 70$
    \item[(3)] $\;\;\;\; N, K, S, A \leq 150$
    \item[(4)] $\;\;\;\; N, K, S, A \leq 1000$
    \item[(5)] $\;\;\;\;$Inga särskilda begränsningar.
\end{itemize}


\fe{Körningsexempel 1}
\begin{verbatim}
N, K, S, A ? 1 1 5 0

Svar: 5
\end{verbatim}
Det finns ett monster, och de möjliga liv den kan ha så att Theodor dödar den med exakt en explosion är \verb+1, 2, 3, 4, 5+.

Se nästa sida för ett till körningsexempel.

\newpage

\fe{Körningsexempel 2}
\begin{verbatim}
N, K, S, A ? 2 2 1 1

Svar: 10
\end{verbatim}
De 10 möjliga fördelningarna av monsterliv så att Theodor dödar båda med exakt två explosioner är:
\begin{verbatim}
  1 3
  1 4
  2 2
  2 3
  2 4
  3 1
  3 2
  3 3
  4 1
  4 2
\end{verbatim}

\end{document}
